
"À l’instar d’autres institutions, les archives audiovisuelles doivent satisfaire à des normes sévères en ce qui concerne la documentation de leurs acquisitions, des communications et d’autres transactions afin de démontrer le sérieux de leur gestion. Du fait de la complexité de leurs fonds, la précision de l’information est ici essentielle. La nature des matériels impose en outre d’établir une documentation précise sur l’utilisation interne des collections." \footcite{zotero-item-772} 

Dans cette philosophie de l'archivistique audiovisuelle publiée par l'UNESCO, Edmonson Ray pose les fondements de la complexité de la gestion et de l'indexation des document audiovisuels en raison de la nature de leur support, une indexation qui doit alors être soumise à une structure définie et claire pour en instaurer la compréhension. En effet, l'auteur définit les documents audiovisuels comme suit : 
"Constituent des documents audiovisuels les œuvres comprenant des images et / ou des sons reproductibles réunis sur un support matériel dont : l’enregistrement, la transmission, la perception et la compréhension exigent le recours à un dispositif technique;le contenu visuel présente une durée linéaire"
\footcite{zotero-item-772}
Ainsi, contrairement aux documents imprimés dont la forme rend accessible le contenu, les supports audiovisuels (films, bandes audio, fichiers vidéo) se caractérisent par une opacité qui nous amène à les qualifier de « boîtes noires » \footcite{zotero-item-582}. Le support physique ne révèle que peu ou pas d'informations sur le contenu du document, qui n'est accessible qu'à sa lecture à l'aide d'un appareil de lecture adéquat, et cela, même une fois le document numérisé. Les institutions patrimoniales préservatrices de fonds audiovisuels sont d'une part, spécifiquement confrontées à la complexité de leur contenu dont l'archiviste doit analyser les images, le son, le montage, la linéarité, etc, et d'autre part, au volume de ces collections, dont la nature du médium ne permet pas d'en décrire le contenu rapidement et précisément. La description du contenu suppose un temps de travail considérable, ainsi qu'un matériel de lecture adéquat. 
L'outil de gestion des collections peut s'imposer comme une solution institutionnelle face à la complexité de ces documents, notamment pour en structurer le contenu, répondant ainsi au premier problème. (deuxième problème : volume : solution = automatisation) En effet, au sein des institutions patrimoniales, la gestion des collections est le plus souvent permise par un outil de gestion que le ministère de la culture définit comme suit : \begin{quote}
	Progiciel spécialisé permettant de gérer des informations structurées et des multimédias sur les objets conservés au sein d’un musée, que ces derniers lui soient affectés ou déposés. C’est un système d’information partagé dont les fonctionnalités correspondent au cycle de vie de l’objet : préparation et gestion de l’entrée dans les collections, prise d’inventaire, documentation, gestion des informations liées à la conservation, récolement, post-récolement, gestion des multimédias, des droits d’auteur, des mouvements, mais aussi diffusion au public (sous forme d'extraction de données ou de catalogue numérique), outil de médiation, etc \footcite{zotero-item-766}
\end{quote}

La collection incarne, selon la même source, un "ensemble non fini d'objets matériels ou immatériels, réunis et classés par un individu ou une institution, en raison de leur valeur scientifique, artistique, esthétique, documentaire, affective, pour leur prix, leur rareté, etc. \footcite{zotero-item-766} Ainsi défini, un outil de gestion des collections permet la réalisation des diverses actions liées au cycle de vie de l'ensemble de ces objets au sein d'une instutition muséale. Toutefois, la possibilité de réaliser ces actions repose nécessairement sur la production "d'informations structurées" exploitables au sein du système. Cette structuration est permise par l'indexation c'est-à-dire « l’opération qui consiste à décrire et à caractériser un document à l’aide de représentations des concepts évoqués dans ce document ». \footcite{zotero-item-767} Il s'agit d'une opération technique mais avant tout intelectuelle. L'indexation permet précisément d'intelligibiliser un objet au sein d'un système de gestion des collections et de le rendre repérable et interrogeable auprès de l'utilisateur. 
La description d'un seul objet existe à plusieurs niveaux : contenu, support, versions. Une indexation simple ne hiérarchise pas ces différents niveaux de description, or, ils sont indispensables à l'appréhension de la gestion de l'objet. C'est pourquoi la modélisation de données permet une hiérarchisation structurée. 
Néanmoins, la structuration du contenu d'un objet par l'extraction de mots-clés ou concepts contrôlés ne suffit pas, dans certains cas, à rendre compte de la complexité de leur contenu. En effet, la nature de certains objets de collections tels que les ouvrages bibliographiques ou les documents audiovisuels, nécessitent que l'indexation de leur contenu s'inscrive au sein d'une modélisation des données. Ainsi, l'outil de gestion des collections peut s'appuyer, pour certaines collections sur un modèle de données, tel que le FRBR (\textit{Functional Requirements for Bibliographic Records}). 

%Dire seulement ce qu'est le FRBR et pourquoi il est utile pour les archive audiovisuelles. 
Le FRBR, issu du milieu des bibliothèques et élaboré par l'IFLA en 1998, puis révisé sous le modèle LRM, est fondé sur une approche entité-relation au sein de laquelle la description de l'objet documentaire en répartie en quatre niveaux : l' \textit{Œuvre} (la création intellectuelle), l' \textit{Expression} (sa réalisation linguistique ou artistique), la \textit{Manifestation} (sa matérialisation éditoriale) et l' \textit{Item} (l'exemplaire physique ou numérique conservé). Ce modèle permet de dissocier le contenu intellectuel d'un objet de son support matériel, dissociation que l'outil de gestion des collections doit concrétiser. Il est adapté, au sein d'un outil de gestion des collections, en fonction de la nature des collections au sein de chaque institution. En effet, si la modélisation FRBR est pensée en premier lieu pour la description d'ouvrage bibliographiques, le modèle est également largement utilisé pour la description des archives et en particulier celle des archives audiovisuelles. Le modèle est alors réadapté en fonction de la nature des objets décrits afin de proposer un mode de description et une structure de données pertinents. Un outil de gestion des collections peut alors proposer une modélisation adaptative à raison des collections de nature multiples préservées par les institutions patrimoniales. 
D'autre part, la modélisation des données audiovisuelles pose un défi théorique tout aussi complexe. Pour s'adapter aux spécificités du médium cinématographique, la Fédération Internationale des Archives du Film (FIAF) a révisé le modèle FRBR original en fusionnant les niveaux \textit{Œuvre} et \textit{Expression} et en introduisant la notion de \textit{Variante} (passant d'une structure WEMI à une structure WVMI). En effet, une œuvre audiovisuelle est indissociable de son expression, mais elle existe à travers une multitude de variantes (montages alternatifs, censures, doublages, restaurations). Si la FIAF propose un cadre théorique, le FRBR reste un modèle extrêmement flexible, réinterprété par chaque institution en fonction de la nature de ses corpus et de ses règles métiers propres.

Réponse au deuxième problème, qu'un outil de gestion des collections peut également permettre : 

Face à l'augmentation croissante des productions de documents, la mise en place de solutions automatiques permettrait de structurer l'information et de pouvoir y donner accès au sein de systèmes de recherche. \footcite{barreaux2017} 

Par ailleurs, les tâches d'indexation permises par un outil de gestion des collections, au sein d'une modélisation de données, répétitives et chronophage, font l'objet de tentatives  d'automatisation à l'aide de l'intelligence artificielle. En effet, face à l'augmentation croissante des productions de documents, la mise en place de solutions automatique permettrait de structurer l'information et de pouvoir y donner accès au sein de systèmes de recherche. \footcite{barreaux2017} Dans le cas de documents textuels ou iconographiques, l'émergence des modèles d'apprentissage automatique au sein du domaine de l'intelligence artificielle, offre des possibilités d'extraction automatique de métadonnées efficaces sur des documents numérisés. Toutefois, certains types de collections telles que les collections audiovisuelles, résistent à la mise en place d'une extraction automatique de leur contenu. En effet, contrairement aux documents imprimés dont la forme rend accessible le contenu, les supports audiovisuels (films, bandes audio, fichiers vidéo) se caractérisent par une opacité qui nous amène à les qualifier de « boîtes noires » \footcite{zotero-item-582}. Le support physique ne révèle que peu ou pas d'informations sur le contenu du document, qui n'est accessible qu'à sa lecture à l'aide d'un appareil de lecture adéquat, et cela, même une fois le document numérisé. 
Ainsi, d'une part, l'indexation automatique de ces fonds supposerait que les modèles d'intelligence artificielle soient capables de lire et de comprendre le langage cinématographique, caractérisé par un enchaînement d'images fixes produisant l'illusion du mouvement, parfois accompagnées de son et d'un montage multiforme. Or, la compréhension automatique de ce langage reste inégale, en particulier lorsqu'elle est confrontée à des fonds d'archives anciens ou dégradés dont la qualité sonore et visuelle est compromise. A priori, l'étude d'un tel type de documents nécessiterait des algorithmes d'intelligence artificielle lourds et couteux, incompatibles avec le milieu patrimoniale. 

Ainsi, automatiser l'indexation des archives audiovisuelles génère une double difficulté : la complexité de lire et de comprendre un contenu visuel et sonore et temporel, et la variabilité d'un modèle de données (FRBR) hautement dépendant du contexte institutionnel.

C'est précisément au coeur de ces enjeux techniques et documentaires que s'inscrit ce mémoire de recherche, réalisé dans le cadre d'un stage de six mois au sein de l'équipe métier de l'entreprise Skinsoft, éditrice de solutions logicielles de gestion des collections, à travers l'étude de son espace de travail dédié aux archive audiovisuelles (S-movies) et la proposition de nouvelles fonctionnalités inscrites dans le cadre d'un projet IA. Cette étude examine les conditions de faisabilité et les modalités d'intégration de modules d'indexation automatique d'images et de sons au sein d'un système de gestion de collections fondé sur une modélisation adaptative du FRBR.

Intégrer un état de l'art : sur l'automatisation de l'indexation des archives audiovisuelles : combinaison de plusieurs algorithmes, pas d'outils concret qui feraient l'enseùble du processus d'indexation. Ce mémoire tente d'articuler un pipeline d'indexation automatique avec un modèle FRBR adaptatif  

Dès lors, nous nous poserons la question suivante : Dans quelle mesure la variabilité institutionnelle et sémantique du modèle de données FRBR adapté aux archives audiovisuelles limite-t-elle la mise en place d'une indexation automatique pour ce type d'archives, et comment l'intelligence hybride permet-elle de dépasser ce verrou ?

Pour répondre à cette interrogation, notre démarche s'articulera en trois temps. 
Nous analyserons d'abord les fondements du modèle FRBR et les modalités théoriques de son adaptation au média audiovisuel, avant d'examiner les opportunités et verrous théoriques posés par l'automatisation de l'indexation.
Nous étudierons ensuite l'implémentation concrète de ces outils au sein du projet d'IA de l'entreprise Skinsoft. Nous analyserons la faisabilité technique du séquençage automatique vidéo (segmentation temporelle) sur des fonds patrimoniaux (comme le corpus du Musée Albert-Kahn) et cartographierons le reversement des métadonnées ainsi produites dans les différents niveaux de la structure FRBR.
Enfin, nous évaluerons les limites, biais et enjeux éthiques liés à l'inadaptation des données d'entraînement d'IA aux fonds d'archives. Nous montrerons alors que les limites du « tout-automatique » invitent à réorienter la conception des logiciels métier vers un paradigme d'**indexation augmentée et d'intelligence hybride**, plaçant l'interaction homme-machine au cœur de la gouvernance de la qualité documentaire.


---------------

	
Cette citation décrit l'indexation comme un travail répétitif d'extraction d'informations directement à partir de la forme et du contenu d'un document. Ce document, ici textuel, contiendrait en lui-même ces informations dès le premier coup d'oeil d'après sa couverture, ses numéros de page, sa table des matières... Ce travail tel qu'il est décrit ici, est aujourd'hui facilement automatisable grâce à des modèles d'apprentissage automatique, capables de lire les documents et en d'extraire les métadonnées pertinentes à indexer pour ensuite les relier à des référentiels. 
	
Or, d'autres types de documents et en particulier les archives audiovisuelles, qui constituent l'objet d'étude de ce mémoire de recherche, résistent à cette automatisation en raison de l'opacité de leur forme, conçue comme une boite noire. En effet, contrairement à la plupart des supports textuels, les archives audiovisuelles ne révèlent, le plus souvent, rien à la lecture de leur contenant. Le contenu apparait au moment de la projection. Ainsi, cette particularité formelle les rend en quelque sorte hermétiques aux pratique d'indexation telles que définies par Rolland-Coul et d'autant plus pour l'automatisation de ces pratiques. \footcite{zotero-item-582}
	
En effet, la mise en place d'une indexation automatisée pour ce type d'archives suppose que les modèles d'intelligence artificielle sachent lire l'enchaînement des images animées et la bande sonore, et en comprenne le langage du montage composé de séquences, de plans, de transitions... Ce langage spécifique n'est pas encore analysé ni compris de manière satisfaisante par les modèles d'intelligence artificielle, surtout dans le cadre de corpus patrimoniaux anciens aux sujets historiques, dont la qualité de l'image et du son est compromise. 
	
À cette difficulté d'extraire automatiquement des métadonnées depuis un contenu audiovisuel, une autre difficulté s'ajoute, celle de leur modélisation. En effet, dans un contexte de gestion patrimoniale des collections audiovisuelles, c'est le modèle FRBR qui est utilisé, issu du milieu des bibliothèques et élaboré par l'IFLA en 1998. IL propose une indexation à quatre niveaux : oeuvre, expression, manifestation, item. Cette indexation par niveaux permet de distinguer les métadonnées relevant de l'oeuvre intellectuelle, des métadonnées physique issues de son contenant. En 2016, la Fédération internationale des archives du film (FIAF) adapte le modèle aux archives audiovisuelles, en fusionnant le niveau Oeuvre/Expression et en ajoutant celui de Variante. Le FRBR n'est plus fondé sur un WEMI mais désormais sur un WVMI. En effet, une oeuvre cinématographique est nécessairement exprimée, mais comporte souvent des variantes : montage différent, changement de personnage, des dialogues. Or, si la FIAF propose un modèle nourri d'exemples et de cas d'application, le FRBR reste un modèle conceptuel flexible, repensé dépendamment du type de corpus indexé et qui diffère ainsi d'une institution à l'autre.
	
	Ainsi, l'automatisation de l'indexation des archives audiovisuelles se heurte à une double difficulté, celle de la compréhension du langage cinématographique et d'un modèle d'indexation de données hautement dépendant du contexte dans lequel il est utilisés. Réalisé dans le cadre d'un stage de six mois au sein de l'entreprise Skinsoft, qui propose des solutions logicielles de gestion des collections, ce mémoire analyse les possibilités d'implémentation d'une indexation automatique au sein d'un espace de travail dédié à la gestion de collections audiovisuelles, fondé sur un modèle de données flexible. 
	
	Dans quelle mesure la variabilité institutionnelle et sémantique du modèle de données FRBR adapté aux archives audiovisuelles limite-t-elle la mise en place d'une indexation automatique pour ce type d'archives ?
	
	Dans un premier temps, nous analyserons le modèle conceptuel du FRBR, un modèle de données issu des bibliothèques, et les modalités de son adaptation aux archive audiovisuelles. Nous verrons quelles questions pose l'implémentation de ce modèle dans le cadre d'une mise en place de l'indexation automatique. Dans un second temps nous proposerons une implémentation concrète d'indexation automatique des archives audiovisuelles au sein de ce modèle de données, dans le cadre du projet IA chez skinsoft. Nous analyserons plus spécifiquement la mise en place de la segmentation vidéo en séquences, une fonctionnalité priorisée dans la mise en place de l'indexation automatique, et sa compatibilité avec le FRBR. Enfin, nous analyserons les biais et limites techniques posés par les lacunes des données d'entraînement des modèles d'intelligence artificielles, difficilement adaptables au contexte patrimonial. Ces limites permettront d'orienter notre réflexion vers une indexation augmentée plutôt qu'automatisée, qui permettrait également d'améliorer la compréhension du FRBR.
	

 

%\chapter{Partie 1 - Indexer les archives audiovisuelles}  
%
%\begin{itemize}
%
%\section{Le FRBR, modèle d’indexation conçu pour les bibliothèques}
%
%
%
%\item {Un modèle entité-relation qui permet la description de concepts reliés entre eux. Avant ce modèle, une notice bibliographique correspondait à un objet physique. Il permet par ailleurs de mieux décrire les objets complexes comme les archives audiovisuelles.}  
%
%\item{La structure WEMI, qui permet de décrire chaque niveau : Oeuvre, expression, manifestation, item. C’est une structure pensée pour des oeuvres textuelles, avec une verticalité limitée }
%
%\item{Les objectifs du FRBR sont les besoins utilisateurs. Il doit pouvoir répondre à quatre tâches utilisateurs : trouver, identifier, sélectionner, obtenir. Afin de remplir ces objectifs, une interface intuitive doit pouvoir accueillir le modèle FRBR.  
%}
%
%\end{itemize}
%
%
%\section{Adapter le FRBR aux archives audiovisuelles}
%
%\begin{itemize}
%
%
%\item{Remodéliser le FRBR sur les caractéristiques complexes du médium cinématographique. L’exemple de S-Movies : un espace de travail pour la gestion des archives audiovisuelles, fondé sur une remodélisation du FRBR, adaptable selon les corpus de données} 
%
%\item{L'inscription de différents types de corpus audiovisuels dans ce modèle repensé du FRBR : la cause d’une expérience utilisateur complexe et non-intuitive }
%
%\item{Le projet d’ajouter sur cette remodélisation du FRBR, une indexation augmentée et automatisée des collections audiovisuelles}
%	
%\end{itemize}
%
%\section{Automatiser l’indexation des collections audiovisuelles basée sur le FRBR} 
%
%\begin{itemize}
%
%
%\item{L’automatisation de l’indexation des archives audiovisuelles, motivée par les caractéristiques physiques de ce type de collection.  En effet, les collections audiovisuelles sont indexées par leur contenu puisque leur support physique est qualifié de “boite noire”. Leur indexation est donc un travail long et fastidieux. } 
%
%\item{Les possibilités d’automatisation de l’indexation sur les collections audiovisuelles : 
%
%Analyse d’objets, (segmentation sémantique) 
%
%Analyse et transcription de texte  
%
%Analyse et tanscription audio  
%
%Reconnaissance vocale et reconnaissance des locuteurs 
%
%Extraction d’entités nommées (lier ces analyses automatiques à des référentiels) }
%
%\item{Séquençage (segmentation temporelle)  : une fonctionnalité spécifique aux archives audiovisuelles puisqu’il s’agit de comprendre les relations entre les plans pour pouvoir adéquatement les séquencer.  
%
%Faire de l’adaptation de cette automatisation au modèle FRBR dédiée aux archives audiovisuelles, un objectif } 
%
%
%\end{itemize}
%
%\chapter{Partie 2 -  Implémenter l’indexation automatique des archives audiovisuelles au sein d’une remodélisation du FRBR} 
%
%
%
%\section{Le projet d’implémentation de l’IA chez Skinsoft (projet Skai), dont l'analyse vidéo est une composante majeure} 
%
%\begin{itemize}
%
%\item{L’objectif du projet est de proposer différents moteurs IA spécialisés dans différents tâches automatisées, entraînés sur des données ouvertes patrimoniales (POP, Joconde, Europeana). Cela permettrait aux modèles d'être adaptés au contexte patrimonial et de produire des métadonnées conformes aux normes bibliographiques. En revanche, ces bases de données contiennent majoritairement des archives textuelles et iconographiques. }
%
%\item{L’orientation donnée est de proposer des modèles force de proposition, et facilitateurs pour l’utilisateur : l’utilisateur doit pouvoir garder le contrôle sur ses données }
%
%\item{Parmi les moteurs IA en cours de développement : un moteur spécialisé pour l’analyse vidéo multimodale, adapté aux collections audiovisuelles. Une des composantes principales de ce moteur, issue d’une demande client, est la segmentation temporelle, c’est-à-dire le séquençage vidéo. L’objectif est, à partir des séquences produites, d’apposer des mots clés sur le contenu de l’image.  }
%
%
%\end{itemize}
%
%
%
%\section{Séquençage et analyse d’un corpus audiovisuel (exemple du corpus de musée Albert Kahn)}
%
%\begin{itemize}
%
%\item{Séquençage des archives audiovisuelles en unités analysables, la première étape complexe d’un pipeline multimodal, dépendante du type de corpus analysé : la détection des coupures permet une classification des segments uis une vérification humaine.}  
%
%\item{Combinaison des approches syntaxique et sémantique pour la segmentation de scènes : réseaux de neurones récurrents (RNN), réseaux de neurones convolutifs (CNN) ou Transformers ? }
%
%\item{L’adaptation de ces outils au corpus Albert Kahn, archives audiovisuelles des années 1920, documentaires et muettes : un séquençage facilement identifiable mais difficilement classifiable d’un point de vue sémantique.} 
%
%\end{itemize}
%
%
%\section{Reversement des métadonnées dans le modèle FRBR : dans quelle mesure les résultats produits peuvent-ils s’inscrire au sein de ce modèle d’indexation ? } 
%
%\begin{itemize}
%
%
%\item{Cartographie du reversement des métadonnées : quels résultats pour quels niveaux du FRBR ? Cette cartographie dépend de l’utilisation du FRBR par l’institution concernée. }
%
%\item{Le reversement adéquat des métadonnées produites sur différents niveaux du FRBR ne peut se produire sans une interface navigable et intuitive entre ces différents niveaux. Une refonte est ainsi nécessaire pour la mise en place d’un reversement automatique. Actuellement, il existe une confusion des liens entre les niveaux, qui fait l’objet d’un projet de refonte dans l’espace de travail S-Movies. } 
%
%\item{La mise en place d’une interface dans laquelle l’utilisateur pourrait valider les données produites, et décider de leur reversement de tel ou tel niveau du FRBR. Le moteur d’IA aurait alors bien ce rôle de “facilitateur” en proposant des solutions, tout en laissant le choix final à l’utilisateur. Cette solution pourrait rendre la compréhenion du FRBR plus claire. } 
%
%
%\end{itemize}
%
%
%\chapter{Partie 3 - Biais, limites et enjeux }
%
%
%
%\section{Les archives audiovisuelles comme données d’entraînement} 
%
%\begin{itemize}
%
%\item{La sous-représentation des archives audiovisuelles anciennes dans les corpus d’entraînement. Il faudrait entraîner le modèle sur les corpus individuels de chaque institution pour pouvoir produire des résultats pertinents. L’inadaptation des corpus d’entraînement entraîne des biais d’indexation.}  
%
%\item{le manque de temps et d’argent dédié à l’entraînement d’un modèle sur de multiples corpus. Quelles solutions ? Fine tuning, transfert learning, données synthétiques. }
%
%\item{Trouver le juste milieu entre un pipeline complexe multimodal qui demande beaucoup de puissance de calcul, et un pipeline léger mais moins performant. Ce choix dépend effectivement du type d’archives à indexer. }
%
%
%\end{itemize}
%
%\section{La difficulté d’adaptation au modèle de données, même avec un entraînement adéquat  }
%
%\begin{itemize}
%
%
%\item{L’utilisation des niveaux du FRBR varie pour chaque institution, puisqu’elle dépend de la nature du corpus audiovisuel. Le modèle est donc interprété par des règles métiers internes et propres à chaque institution. Il faudrait ainsi entraîner un modèle IA sur ces règles métier liées à l’utilisation du FRBR. Par exemple, la collection du musée Albert Kahn s’adapte mal au FRBR puisqu’il s’agit d’archives documentaires, sans réelle diffusion. Un reversement automatique suppose que le modèle d’IA soit entraîné sur le contexte spécifique de chaque institution. }
%
%\item{S’impose alors la mise en place d'une indexation semi automatisée où l’utilisateur choisit, décide pour respecter la cohérence de l’usage du modèle de données. }
%
%
%\end{itemize}
%
%\section{Vers une indexation augmentée plutôt qu’automatisée }
%
%\begin{itemize}
%
%\item{L’interactivité et le dialogue entre le modèle et l’utilisateur comme réponse aux biais et aux manques des données d’entraînement. Il faudrait rendre le processus transparent comme le préconise Skinsoft dans son projet Skai, en rendant accessible les données sur lequelles le modèle a été entraîné, la réflexion effectuée... }
%
%\item{La transparence comme éthique de la transmission du savoir au public : la diffusion des connaissances est facilitée et augmentée par le modèle d’IA. La mise en place d’une interface de dialogue entre le modèle et l’utilisateur en est la condition essentielle }
%
%\end{itemize}	



